\documentclass[11pt,a4paper]{article}
\usepackage[T1]{fontenc}
\usepackage{amsmath,amssymb,amsthm}
\usepackage{array}
\usepackage{geometry}
\geometry{margin=2.5cm}

\DeclareMathOperator{\Var}{Var}
\newtheorem{theorem}{Theorem}
\newtheorem{proposition}[theorem]{Proposition}
\newtheorem{lemma}[theorem]{Lemma}

\title{Theoretical Research Report: EABC-Proof-Theory (Version 19)}
\author{Research Collaboration}
\date{2026}

\begin{document}
\maketitle

\begin{abstract}
This document establishes formal-logical and algebraic components of the
Euler-Adic-Borel-Cantor (EABC) research program. We formulate the orthogonal
Pythagorean ANOVA decomposition of the Hurwitz--Hoffbauer defect, derive its
exact representation as an unnormalized window-integrated algebraic two-point
functional, prove that the uniform asymptotic frequency set over the four-symbol
EABC alphabet is $\mathbf{\Pi}^0_3$-complete, and analyze a specified
eight-state isotropic reward-controlled Markov kernel. For this kernel the
limiting entropy is positive, so deterministic zero-entropy crystallization is
excluded.
\end{abstract}

\section{Pythagorean Orthogonal Decomposition and Variance Splitting}

Let $\eta_{\min}:=(1,1,1,1)$ denote the minimal EABC reference pattern,
and let $P=(h_E,h_A,h_B,h_C)$ be the observed Hoffbauer bracket. Define
\[
\delta:=P-\eta_{\min}
=(\delta_E,\delta_A,\delta_B,\delta_C),
\qquad
\delta_X:=h_X-1,
\]
and
\[
K:=\langle\delta,\mathbf 1_4\rangle
=\sum_{X\in\{E,A,B,C\}}\delta_X,
\qquad
\mathbf 1_4:=(1,1,1,1).
\]

Since $\|\mathbf 1_4\|_2^2=4$, the space $\mathbb R^4$ decomposes as
\[
\mathbb R^4
=\operatorname{span}\{\mathbf 1_4\}
\mathbin{\overset{\perp}{\oplus}}
\mathbf 1_4^\perp,
\qquad
\delta
=\underbrace{\frac K4\mathbf 1_4}_{\delta_\parallel}
+\underbrace{\left(\delta-\frac K4\mathbf 1_4\right)}_{\delta_\perp}.
\]
By construction, $\langle\delta_\perp,\mathbf 1_4\rangle=0$.

Define the unminimized direct defect by
\[
D_{\mathrm{HH}}^0(P)
:=\|\delta\|_2^2
=\sum_{X\in\{E,A,B,C\}}\delta_X^2.
\]
The Pythagorean theorem gives
\[
\boxed{
D_{\mathrm{HH}}^0(P)
=\underbrace{\frac{K^2}{4}}_{\text{radial quadratic contribution}}
+\underbrace{D_{\mathrm{HH}}^\perp(P)}_{
\text{internal EABC asymmetry}}
}
\]
with
\[
\begin{aligned}
D_{\mathrm{HH}}^\perp(P)
&:=\|\delta_\perp\|_2^2
=\sum_X\left(\delta_X-\frac K4\right)^2\\
&=\sum_X\delta_X^2-\frac14\left(\sum_X\delta_X\right)^2
=4\,\Var_{\mathrm{pop}}(\delta_E,\delta_A,\delta_B,\delta_C).
\end{aligned}
\]

\subsection{Representation as an Integrated Two-Point Functional}

Let $I$ be a finite window of EABC prime classes. Put
\[
N:=|I|,
\qquad
h_X:=\#\{t\in I:c_t=X\}.
\]
Then $\sum_Xh_X=N$ and $K=\sum_X(h_X-1)=N-4$. For the centered
simplex embedding
\[
v(c_t):=e_{c_t}-\frac14\mathbf 1_4\in\mathbf 1_4^\perp
\]
one has
\[
\sum_{t\in I}v(c_t)
=\left(h_E-\frac N4,h_A-\frac N4,h_B-\frac N4,h_C-\frac N4\right)
=\delta_\perp.
\]
Consequently,
\[
\boxed{
\begin{aligned}
D_{\mathrm{HH}}^\perp(I)
&=\left\|\sum_{t\in I}v(c_t)\right\|_2^2\\
&=\sum_{r,s\in I}\langle v(c_r),v(c_s)\rangle\\
&=\sum_{X\in\{E,A,B,C\}}h_X^2-\frac{N^2}{4}.
\end{aligned}}
\]
Thus $D_{\mathrm{HH}}^\perp$ is an unnormalized, window-integrated
algebraic two-point functional. No physical or thermodynamic assumptions are
required for this identity.

\subsection{Range of the Relative Anisotropy Index}

For $D_{\mathrm{HH}}^0(P)>0$, define
\[
R_\perp(P)
:=\frac{D_{\mathrm{HH}}^\perp(P)}{D_{\mathrm{HH}}^0(P)}
=1-\frac{K^2}{4D_{\mathrm{HH}}^0(P)}.
\]
It is invariant under nonzero rescaling $\delta\mapsto\lambda\delta$.

\begin{enumerate}
\item In the ambient defect space $\mathbb R^4$, Cauchy--Schwarz gives
\[
0\le R_\perp\le1,
\]
with
\[
R_\perp=0\iff\delta\in\operatorname{span}\{\mathbf 1_4\},
\qquad
R_\perp=1\iff K=0\ \text{and}\ \delta\ne0.
\]

\item In the positive excess cone $\delta_X\ge0$, nonnegativity of the
cross terms and Cauchy--Schwarz give
\[
D_{\mathrm{HH}}^0(P)\le K^2\le4D_{\mathrm{HH}}^0(P),
\qquad
\boxed{0\le R_\perp\le\frac34}.
\]
Here $R_\perp=0$ if and only if all four components of $\delta$ are
equal. For a nonzero integer-valued radial defect this implies $4\mid K$.
The equality $R_\perp=3/4$ holds if and only if exactly one component is
nonzero; up to permutation, $\delta=(K,0,0,0)$ with $K>0$.
\end{enumerate}

\section{Descriptive Complexity and Topological Classification}

Let $\mathcal C:=\{E,A,B,C\}$ and let
$Y:=\mathcal C^{\mathbb N}$ carry the product topology. Define
\[
A_{\mathrm{freq}}
:=\left\{\omega\in Y:
\lim_{n\to\infty}\frac1n
\sum_{i=1}^n\mathbf 1_{\{X\}}(\omega_i)=\frac14
\text{ for every }X\in\mathcal C\right\}.
\]

We use the standard prototype
\[
\mathcal Z
:=\left\{\alpha\in2^{\mathbb N}:
\lim_{n\to\infty}\frac1n\sum_{i=1}^n\alpha_i=0\right\},
\]
which is $\mathbf{\Pi}^0_3$-complete; see \cite{KiLinton1994}.

\begin{theorem}
The uniform asymptotic frequency set $A_{\mathrm{freq}}$ is a
$\mathbf{\Pi}^0_3$-complete (equivalently, $F_{\sigma\delta}$) subset of
$\mathcal C^{\mathbb N}$.
\end{theorem}

\begin{proof}
For the upper bound, for $X\in\mathcal C$, $m\ge1$, and $N\ge1$, let
\[
F_{X,m,N}
:=\bigcap_{n\ge N}
\left\{\omega:\left|
\frac1n\sum_{i=1}^n\mathbf 1_{\{X\}}(\omega_i)-\frac14
\right|\le\frac1m\right\}.
\]
Each set inside the intersection depends on a finite prefix and is clopen;
hence $F_{X,m,N}$ is closed. Moreover,
\[
A_{\mathrm{freq}}
=\bigcap_{X\in\mathcal C}\bigcap_{m=1}^\infty
\bigcup_{N=1}^\infty F_{X,m,N},
\]
so $A_{\mathrm{freq}}\in\mathbf{\Pi}^0_3$.

For hardness, define a continuous map
$f:2^{\mathbb N}\to\mathcal C^{\mathbb N}$ using blocks of constant length
four. Let
\[
H:=(E,A,B,C),
\qquad
S:=(E,E,E,E),
\]
and set
\[
f(\alpha):=B_1B_2B_3\cdots,
\qquad
B_k:=\begin{cases}
H,&\alpha_k=0,\\
S,&\alpha_k=1.
\end{cases}
\]
Every output prefix depends on only finitely many input bits, so $f$ is
continuous. Write $d_n:=n^{-1}\sum_{k=1}^n\alpha_k$. At the end of the
$n$th block,
\[
\operatorname{freq}_{4n}(X;f(\alpha))
=\frac{1-d_n}{4}
\quad(X\in\{A,B,C\}),
\]
whereas
\[
\operatorname{freq}_{4n}(E;f(\alpha))
=\frac14+\frac34d_n.
\]
The discrepancy between an arbitrary prefix and the preceding block endpoint
is at most $4/N$, which tends to zero. Hence
\[
\alpha\in\mathcal Z
\iff d_n\to0
\iff f(\alpha)\in A_{\mathrm{freq}}.
\]
Thus $\mathcal Z\le_{\mathrm W}A_{\mathrm{freq}}$, proving
$\mathbf{\Pi}^0_3$-hardness and therefore completeness.
\end{proof}

\section{Markov-Kernel Analysis and Entropy Limits}

Let $S_8:=\mathcal C\times\{+,-\}$, so $|S_8|=8$. This is a specified
model state space; no arithmetic identification beyond this product definition
is asserted. For fixed $T_{\mathrm{eff}}>0$ and scale parameter $\rho>0$, put
\[
\beta_\rho:=\frac{\rho^2}{T_{\mathrm{eff}}},
\qquad
E_{ij}(\rho):=-\rho^2\kappa_{ij},
\qquad
\kappa_{ii}=0,
\qquad
\kappa_{ij}=1\quad(i\ne j).
\]
The associated Gibbs kernel is
\[
T_{ij}(\rho)
=\frac{\exp(\beta_\rho\kappa_{ij})}
{\sum_{k=1}^8\exp(\beta_\rho\kappa_{ik})}.
\]
As $\rho\to\infty$,
\[
T(\rho)\longrightarrow T_\infty:=\frac17(J_8-I_8).
\]

\subsection{Entropy and Information Limits}

The formal high-temperature limit $\beta_\rho\to0$ has row entropy
\[
H_{\mathrm{high}}=\log_2 8=3\ \mathrm{bits}.
\]
In the low-effective-temperature limit $\beta_\rho\to\infty$,
\[
H_\infty=\log_2 7\approx2.80735\ \mathrm{bits},
\qquad
\Delta H=\log_2\frac87\approx0.192645\ \mathrm{bits}.
\]
Under the uniform stationary distribution of the limiting chain,
\[
I(X_n;X_{n+1})=\log_2\frac87.
\]
This mutual information records exactly the exclusion of the current state
from the seven admissible successors.

\subsection{Spectral Characterization}

The spectrum is
\[
\operatorname{spec}(T_\infty)
=\left\{1,\left(-\frac17\right)^{[7]}\right\}.
\]
With $\Pi:=J_8/8$,
\[
T_\infty^n
=\Pi+\left(-\frac17\right)^n(I_8-\Pi).
\]
Thus every initial distribution converges geometrically to the uniform
stationary distribution. The limiting kernel describes isotropic
self-avoidance, not deterministic crystallization.

\section{Epistemic Status Summary}

\begin{center}
\begin{tabular}{>{\raggedright\arraybackslash}p{0.62\linewidth}|
                >{\raggedright\arraybackslash}p{0.24\linewidth}}
\hline
\textbf{Component} & \textbf{Status}\\
\hline
Pythagorean defect decomposition & $[\mathrm T]$ Proved\\
Integrated two-point functional identity & $[\mathrm T]$ Proved\\
Anisotropy bounds in the stated domains & $[\mathrm T]$ Proved\\
Uniform frequency set is $\mathbf{\Pi}^0_3$-complete
& $[\mathrm T]$ Proved, using the cited prototype theorem\\
Specified Markov limit $T_\infty=(J_8-I_8)/7$ & $[\mathrm T]$ Proved\\
Deterministic zero-entropy limit for this kernel & $[\mathrm F]$ Refuted\\
Application of the eight-state kernel to prime dynamics & $[\mathrm H]$ Open\\
\hline
\end{tabular}
\end{center}

\begin{thebibliography}{9}
\bibitem{KiLinton1994}
H.~Ki and T.~Linton,
\emph{Normal numbers and subsets of $\mathbb N$ with given densities},
Fundamenta Mathematicae \textbf{144} (1994), 163--179.
\end{thebibliography}

\end{document}
